% =============================================================================
% 3.1 Présentation de la plateforme
% =============================================================================
\section{Présentation de la plateforme développée}
\label{sec:ch3-plateforme}

La plateforme implémente la méthodologie du chapitre~\ref{chap:methodologie}.
Elle permet l'ingestion de documents hétérogènes, l'indexation conjointe
lexicale et vectorielle, la recherche hybride, la génération contrainte
de sections, la revue humaine et l'export d'un livrable versionné.

\subsection{Interface utilisateur}

L'interface s'organise en quatre espaces, qui correspondent aux quatre
capacités identifiées en méthodologie.

\begin{itemize}
  \item un espace de dépôt et de gestion des documents, avec indication du
        statut d'indexation (en attente, en cours, prêt, en échec)~;
  \item un espace de conversation, point d'entrée unique pour les questions,
        les demandes de génération et les exports~;
  \item un espace de revue, où chaque section générée est présentée avec
        ses citations, ses exigences couvertes et ses abstentions~;
  \item un espace d'export, qui n'assemble que les sections approuvées.
\end{itemize}

\subsection{Périmètre de la version évaluée}

La version évaluée dans ce chapitre couvre les phases 1 à 4 du phasage.
Le graphe de connaissances n'est pas activé. Les mesures portent donc sur
le RAG hybride, le découpage parent-enfant, le ré-ordonnancement et la
revue humaine, et non sur des questions relationnelles multi-sauts.
